% 3 Data and methods (all; first draft: Aryan). Sources: deck slides 10-14; scripts in presentation/analysis and
% report/analysis. Parameter values: presentation/analysis/water_params.py (claims.csv rows O2c).
\section{Data and methods}\label{sec:methods}

\subsection{Datasets}
\Tab{tab:data} lists the sources actually used. The backbone is Epoch AI's database of AI data centres
\citep{epoch2026dc}: 96 large AI sites with their current facility power and build timelines to the end of 2030. It is
the only open dataset that gives power and timelines site by site. Its coverage is the main limit on what follows
(Section~\ref{sec:limits}): it tracks AI sites only, not all data centres, and most of them are in the United States.

\begin{table}[t]\centering
\caption{Datasets used, by theme. All are public; each number in the report traces to one of them or to our scripts.}
\label{tab:data}
\footnotesize\renewcommand{\arraystretch}{1.08}
\begin{tabular}{@{}P{1.75cm}P{5.85cm}@{}}
\toprule
Theme & Sources \\
\midrule
Sites and capacity & Epoch AI data centres: 96 sites, facility MW today and planned to 2030 \citep{epoch2026dc};
geocoded with OpenStreetMap Nominatim \citep{osm2026} \\
Electricity, efficiency & IEA global and regional, 2024--2030 \citep{iea2025}; LBNL US series 2014--2023
\citep{shehabi2024}; global estimates 2010--2018 \citep{masanet2020}; Uptime PUE survey 2007--2025
\citep{uptime2025} \\
Growth & World primary energy 1950--2019 \citep{owid2026energy}; training compute of 540 notable AI models
\citep{epoch2026models} \\
Carbon & Ember grid carbon intensity by country, 2024 \citep{ember2026} \\
Water & WRI Aqueduct 4.0: baseline annual and monthly stress, 2030 and 2050 projections \citep{kuzma2023}; WUE, PUE
and grid water-intensity ranges \citep{devriesgao2026}; city records and news reports (case studies) \\
Population & GHSL GHS-POP 2025, 30 arc-second grid \citep{schiavina2023} \\
Materials, waste & USGS Mineral Commodity Summaries 2026 \citep{usgs2026}; Global E-waste Monitor 2024
\citep{balde2024} \\
Boundaries & Planetary Health Check 2025 \citep{sakschewski2025}; Earth3 thresholds \citep{collste2018}; Great
Acceleration framing \citep{steffen2015ga} \\
Society & Data Center Watch Q1 and Q2 2026 (US) \citep{dcw2026,dcw2026q2}; EJAtlas grievance categories \citep{temper2015} \\
Siting record & Public statements, official documents and press on the Visakhapatnam hub; Google's 2026 Environmental Report
\citep{google2026env}; NITI Aayog framework \citep{niti2026}; CEEW study \citep{ceew2026} (Section~\ref{sec:siting}) \\
\bottomrule
\end{tabular}
\end{table}

\paragraph{Changes from the plan} The plan proposed a wider net. We used Epoch AI sites only, with no all-data-centre
layer and no footprint-based size model; two case studies instead of three or four clusters, plus a siting audit of the Visakhapatnam hub (Section~\ref{sec:siting}), which was not in the plan; country population from
the GHSL grid rather than UN WPP for the location quotients; no Gini per internet user and no dispute catalogue of our
own; and world primary energy as the only Great Acceleration series.

\subsection{Analytical and statistical methods}

\paragraph{1. Spatial overlay} Each Epoch site was geocoded from its address (all 96 matched, with no country or
state mismatches after restricting queries by country code) and joined to its Aqueduct basin. City-level accuracy is
enough: basins are tens to hundreds of kilometres across. For every site we read the baseline water-stress category,
the 2030 and 2050 projections (business as usual and pessimistic) and the number of months in which withdrawals exceed
supply. The key output is the share of capacity, weighted by MW, in high or extremely high stress, now and for the
planned 2030 capacity. People near the sites were counted from GHSL cells whose centres lie within 25 and 50\,km;
overlapping circles are merged, so nobody is counted twice. The same cells, aggregated by country, are used for the
country table (method 2).

\paragraph{2. Inequality} For each country $i$ with capacity $C_i$ and population $P_i$, the location quotient
\begin{equation}
LQ_i = \frac{C_i / C}{P_i / P}
\end{equation}
is above 1 when a country hosts more than its population share. Countries are then sorted by capacity per person to
draw a Lorenz curve, with cumulative population share $x_k$ and cumulative capacity share $y_k$, and
\begin{equation}
G = 1 - \textstyle\sum_k (x_k - x_{k-1})(y_k + y_{k-1}).
\end{equation}
Countries with no tracked site are pooled into one zero-capacity group. Because they all have exactly zero capacity,
pooling them leaves $G$ unchanged, so this is the exact country-level Gini for the Epoch sample. Country and world
population are sums of GHSL 2025 cells (8.19 billion people), so numerators and denominators come from the same grid.

\paragraph{3. Water-use range} Water is estimated from electricity. With facility power $M$ (MW) and load factor
$L$,
\begin{equation}
E_{\text{fac}} = M L \times 8{,}760\,\text{h},\qquad E_{\text{IT}} = E_{\text{fac}} / \text{PUE},
\end{equation}
\begin{equation}
W_{\text{direct}} = E_{\text{IT}}\times\text{WUE},\qquad W_{\text{indirect}} = E_{\text{fac}}\times WI_{\text{grid}}.
\end{equation}
The low case uses PUE 1.10, WUE 0.56\,L/kWh and $WI_{\text{grid}}$ 1.04\,L/kWh; the high case PUE 1.43, WUE 1.04
and $WI_{\text{grid}}$ 4.52, all from \citet{devriesgao2026}, who compiles IEA, LBNL and company values. The load factor
(70\% low, 90\% high) is \emph{our assumption}: AI clusters run close to continuously, but no site reports its average
draw. Both terms are \emph{consumption}. The low and high cases combine the extremes of each range, so they are bounds,
not a confidence interval.

\paragraph{4. Standardised growth} Each series is fitted as $\ln y_t = a + bt$ by ordinary least squares, giving the
annual growth rate and doubling time
\begin{equation}
r = e^{b} - 1,\qquad T_2 = \ln 2 / b.
\end{equation}
Intervals (95\%) are reported where a yearly series exists; observed and projected years are kept separate. For the
report we added fit diagnostics: $R^2$, the Durbin--Watson statistic for serial correlation in the residuals
\citep{durbin1950}, and a Newey--West standard error for $b$ \citep{newey1987} with
$\lfloor 4(n/100)^{2/9}\rfloor$ lags. Serial correlation leaves the slope unbiased but makes the ordinary interval too
narrow. To test where a trend breaks, a continuous two-segment (hinge) model
$\ln y_t = a + bt + c\,(t-\tau)_+$ is fitted for every candidate year $\tau$ and the $\tau$ with the smallest residual sum
of squares is kept. The change in the Bayesian information criterion \citep{schwarz1978} between one and two segments
(the break year counted as a parameter) says whether the break is worth having: the more negative, the stronger the case.

\paragraph{5. Case studies and synthesis} Two contrasting cases were built from official records and news reports:
Google's data centre in Mesa, Arizona, which is on our map, and the planned Google--AdaniConneX hub in Visakhapatnam,
India, which is not yet in the Epoch data. For each we record the water source and its distance, the local reaction,
and our own estimate. All results are then placed in one matrix of planetary boundary against scale of impact, with each
cell coded \textbf{Q} (quantified in this project), \textbf{D} (documented pathway, not quantified) or \textbf{N} (not
assessed).

\paragraph{6. Siting audit} For the Visakhapatnam hub we read what its backers and the authorities have said about the choice
of site, and list the instruments that could have governed it. We then place the hub's Aqueduct basin class among those of the 96
tracked sites, weighted by planned MW, and rerun the water model of method 3 for 1\,GW under two cooling designs: wet cooling,
and air cooling that needs 10 or 20\% more energy, so that indirect water rises in proportion. The 10\% is Google's own average
figure; the 20\% is a test for a hot climate. Only the arithmetic is ours (\texttt{analysis/fig\_siting.py}); the reading of the record
rests on public sources, which we cite, and is limited to what they say.
